The theorem of Papadakis and Petrotou applies to simplicial homology spheres $\Delta$ over a field $k$ of characteristic $2$. The Stanley-Reisner ring $\cA(\Delta)$ and the cohomology ring $H(\Delta) = \cA(\Delta)/(\theta_1,\ldots,\theta_n)$ are defined over a larger field $K=k(\a{i,j})$ of rational functions in the variables $\a{i,j}$. The variables $\a{i,j}$ here are the coefficients of the linear parameters $\theta_1,\ldots, \theta_n$ in the definition of $H(\Delta)$. 

\begin{theorem}[Papadakis, Petrotou] \label{thm-PP}
Let $\Delta$ be a simplicial homology sphere of dimension $n-1=2m-1$ over a field $k$ of characteristic $2$. Let $H(\Delta)$ be defined over the field of rational functions $K=k(\a{i,j})$. Then the quadratic form defined on the middle degree cohomology $H^m(\Delta)$ by multiplication 
\[ Q(g) = g^2 \in H^{n}(\Delta) \isom K\]
 is anisotropic.
\end{theorem}

Papadakis and Petrotou used Theorem~\ref{thm-PP}  to prove the Hard Lefschetz theorem for all simplicial homology spheres in characteristic $2$,  in both even and odd dimensions. The Hard Lefschetz theorem then implies the g-conjecture for such spheres. 

Our first result in this article is an explicit description of the quadratic form $Q$ that holds in any characteristic (Theorem~\ref{thm-Q} below). We use this description to prove in Theorem~\ref{thm-PPconj} a conjecture stated in \cite{PapadakisPetrotou} that generalizes the main ingredient in the proof of Theorem~\ref{thm-PP}. As an application of the conjecture, we prove anisotropy in all degrees $m\leq n/2$.

One can define the Hodge-Riemann type quadratic form in any degree. Let
\[ l = x_1 + x_2 + \cdots + x_N \in H^1(\Delta),\]
where $\Delta$ has $N$ vertices and $x_1,\ldots, x_N$ are the corresponding variables in the Stanley-Reisner ring. 
The quadratic form $Q_l$ on $H^m(\Delta)$ for $m\leq n/2$ is defined by
\[ Q_l(g) = l^{n-2m} g^2 \in H^n(\Delta) \isom K.\]
%This form is defined for both even and odd $n$. 

\begin{theorem} \label{thm-lef}
Let $\Delta$ be a simplicial homology sphere of dimension $n-1$ over a field $k$ of characteristic $2$, and let $H(\Delta)$ be defined over the field of rational functions $K=k(\a{i,j})$. Then the quadratic form $Q_l$ is anisotropic on $H^m(\Delta)$ for any $m\leq n/2$. 
\end{theorem}

Theorem~\ref{thm-lef} can be deduced from Theorem~\ref{thm-PP} using induction on the dimension of the sphere and Hard Lefschetz theorem \cite{PapadakisPetrotou}. However, Theorem~\ref{thm-lef} is also a simple application of the conjecture in \cite{PapadakisPetrotou}. 
%The conjecture is indeed well suited for proving all types of anisotropy results in characteristic $2$. 
Note also that Theorem~\ref{thm-lef} directly implies the Hard Lefschetz theorem, which is equivalent to the form $Q_l$ being nondegenerate on $H^m(\Delta)$ for $m \leq n/2$.
